% Checkpoint34 proposed source. All earlier sources are unchanged.
\statement{A.109}{Proved}{The actual derived cyclic cohomology has two points, and the original CM sign is its specified transporter obstruction}
Use the original
\[
 B=C_2(s)\amalg C_3(r),\qquad F=\ker(B\longrightarrow S_3),
 \qquad x=(sr)^{-2},\quad y=sxs,
\]
where \(F\) is the full free profinite group and \(F'\) its closed
derived subgroup. Put \(\rho=\operatorname{Ad}(r)|_F\), and retain
the actual signed map
\[
 q_3(x)=X=\begin{pmatrix}1&2\\0&1\end{pmatrix},\qquad
 q_3(y)=Y=\begin{pmatrix}1&0\\1&1\end{pmatrix}
 \quad\text{over }\mathbf F_3.
\]
Write \(C=XYX^{-1}=J\bmod3\), with the integral
\(J=\left(\begin{smallmatrix}0&-1\\1&-1\end{smallmatrix}\right)\),
and put \(K'_3=\ker(q_3|_{F'})\). Then
\[
 (F')^\rho=1,\qquad H^1(C_3,F')=\{*\},\qquad
 \{I,-I\}\ \xrightarrow{\ \sim\ }\ H^1(C_3,K'_3).
 \tag{A.109.1}
\]
The last map is a pointed homeomorphism, given by
\[
 a\longmapsto[g^{-1}\rho(g)],\qquad
 g\in F',\quad q_3(g)=a.
 \tag{A.109.2}
\]
Here cocycles satisfy \(z\rho(z)\rho^2(z)=1\), and equivalence is
\(z\sim b^{-1}z\rho(b)\). No group structure on \(H^1\) is assumed.

For the actual negative \(\alpha=(-1,\phi)\in\widehat{\mathrm{GT}}_0\),
retain the original LM7\_PUSH~E.61/E.65 frames of A.105. Thus
\[
 \phi=\rho(A_r)A_r^{-1},\qquad
 q_3(A_r)=\epsilon_A X^\nu C^{-\nu},\qquad A_r\in F'.
\]
For \(\epsilon\in\{1,-1\}\), choose any \(t\in F'\) with
\(q_3(t)=a_\epsilon:=\epsilon X^\nu C^{-\nu}\). The class
\[
 \operatorname{ob}(a_\epsilon)
   =[t^{-1}\phi^{-1}\rho(t)]\in H^1(C_3,K'_3)
 \tag{A.109.3}
\]
is independent of the lift and corresponds in (A.109.1) to
\(\epsilon_A\epsilon I\). Consequently
\[
 \operatorname{ob}(a_\epsilon)=*
       \quad\Longleftrightarrow\quad q_3(A_r)=a_\epsilon.
 \tag{A.109.4}
\]
For full GT, with the actual \(d=1+2k\), the label of
\(\operatorname{ob}(a_{d\bmod3})\) is
\((-1)^{\Delta(\alpha)}I\). The choice of this point remains open.
Every actual continuous \(\rho\)-equivariant finite \(p\)-group
quotient of \(K'_3\), for \(p\ne3\), sends both classes to the
basepoint. This last assertion concerns the transporter equation
alone, not a target retaining additional full-pentagon data.
\begin{proof}
\emph{A.109.1. The original cyclic words and both torsion branches.}
The based words give
\[
 \begin{gathered}
 x=r^{-1}sr^{-1}s,\qquad y=sr^{-1}sr^{-1},\qquad
 sr^{-1}s=yr,\\
 \rho(x)=y,\quad \rho(y)=(xy)^{-1},\quad
 \rho((xy)^{-1})=x,\\
 \rho_{\rm ab}=\begin{pmatrix}0&-1\\1&-1\end{pmatrix},
 \qquad
 I-\rho_{\rm ab}=\begin{pmatrix}1&1\\-1&2\end{pmatrix}.
 \end{gathered}
 \tag{A.109.5}
\]
The finite-factor torsion and centralizer theorems
\cite[Theorems A' and B']{A_HR1985} imply that an element of order
three in \(B\) is conjugate to \(r\) or \(r^{-1}\), and that
\(C_B(r)=C_3(r)\). The original \(S_3\) quotient is injective on
this factor. Hence \(F^\rho=(F')^\rho=1\).

Take \(h\in F'\) with \(h\rho(h)\rho^2(h)=1\). The element \(hr\)
has order three and has \(S_3\)-image \(\bar r\).
If \(hr=grg^{-1}\), then \(\bar g\in\langle\bar r\rangle\).
Right multiplication of \(g\) by a suitable power of \(r\)
does not change its conjugate and makes it a pure \(f\in F\).
Thus \(h=f\rho(f)^{-1}\). Abelianization gives
\[
 (I-\rho_{\rm ab})f_{\rm ab}=0.
 \]
Its determinant is three, and multiplication by three is injective
on \(\widehat{\mathbf Z}\). Therefore \(f\in F'\), so
\(h=(f^{-1})^{-1}\rho(f^{-1})\) is a coboundary in \(F'\).

If \(hr=gr^{-1}g^{-1}\), the \(S_3\)-image of \(g\) is in
\(\bar s\langle\bar r\rangle\). Right multiplication by a power of
\(r\) normalizes it to \(\bar s\); write \(g=fs\), \(f\in F\).
The literal identity in (A.109.5) gives
\[
 h=f y\rho(f)^{-1}.
 \]
Writing \(f_{\rm ab}=(a,b)\), its vanishing abelianization requires
\[
 a+b=0,\qquad -a+2b=-1,\qquad\text{hence }3b=-1.
 \]
This is impossible in the three-adic coordinate of
\(\widehat{\mathbf Z}\). The second branch is excluded. This proves
\(H^1(C_3,F')=\{*\}\) using only the actual finite-factor theorem.

\emph{A.109.2. The actual derived quotient and the complete two-point map.}
The quaternion calculation in A.105 gives
\[
 q_3(F')=Q_8,\qquad C_{Q_8}(C)=\{I,-I\}.
 \tag{A.109.6}
\]
The actual cyclic substitution on all \(F\) differs from
\(\operatorname{Ad}(C)\) by a central sign character. On \(F'\)
that character is one, so \(q_3\rho=\operatorname{Ad}(C)q_3\)
and \(K'_3\) is stable. This derived restriction is essential;
no invariance of \(\ker q_3\) on all \(F\) is assumed.

For \(a=\pm I\), a lift \(g\in F'\) exists by (A.109.6).
Its defect \(g^{-1}\rho(g)\) is a \(K'_3\)-cocycle.
Changing the lift to \(gn\), \(n\in K'_3\), gives the cohomologous
defect \(n^{-1}g^{-1}\rho(g)\rho(n)\).
Every cocycle in \(K'_3\) is a coboundary in \(F'\) by the first
part; the image of its conjugating element must be fixed by
\(\operatorname{Ad}(C)\), hence is \(I\) or \(-I\).
This proves surjectivity.

If the two defects become equal after \(K'_3\)-conjugation, absorb
that conjugator in one lift to obtain
\[
 g_1^{-1}\rho(g_1)=g_2^{-1}\rho(g_2).
 \]
Then \(g_2g_1^{-1}\in(F')^\rho=1\), so their original finite
labels were equal. This proves injectivity directly.
The cocycle set is compact and its compact-group orbit relation
is closed; \(H^1\) is therefore compact Hausdorff. The finite
pointed bijection is a homeomorphism.

\emph{A.109.3. The original frame and its exact obstruction.}
A.105 proves \(A_r\in F'\): the relation
\(\phi=\rho(A_r)A_r^{-1}\) and \(\phi\in F'\) give
\((\rho_{\rm ab}-I)(A_r)_{\rm ab}=0\), and the same determinant
argument applies. It also proves the two actual finite candidates
\(a_\epsilon\). Both lie in \(Q_8\): the \(C_3\) abelianization
sends \(X,C\) to one and \(-I\) to zero. They satisfy
\[
 a_\epsilon C a_\epsilon^{-1}=q_3(\phi)^{-1}C.
 \]
Thus (A.109.3) is a cocycle in the correctly typed \(K'_3\).
The norm follows either by expansion or from
\((\phi^{-1}r)^3=1\). A lift change \(t\mapsto tn\) gives the
usual \(K'_3\)-cohomology equivalence.

Put \(b=A_r^{-1}t\). In the actual order,
\[
 t^{-1}\phi^{-1}\rho(t)=b^{-1}\rho(b),\qquad
 q_3(b)=\epsilon_A\epsilon I.
 \tag{A.109.7}
\]
This identifies the class in (A.109.1), and proves (A.109.4).
Equivalently, triviality gives \(tn\) satisfying the original
frame equation; uniqueness follows from \(C_B(r)\cap F=1\).
For full GT, A.105 gives
\(\epsilon_u=d\bmod3\) and
\((-1)^{\Delta(\alpha)}=\epsilon_A\epsilon_u\), proving the
stated original CM label. It assigns no value to that label.

\emph{A.109.4. Why coprime-\(p\) transporter tests cannot select it.}
For a finite \(p\)-group \(P\) with an order-three action and
\(p\ne3\), we prove \(H^1(C_3,P)=\{*\}\) by induction on its order.
For nontrivial \(P\), the nonzero characteristic subgroup
\[
 Z=\Omega_1(Z(P))
 \]
is elementary abelian, central and stable. Induction modulo \(Z\)
allows conjugation of the cocycle into \(Z\). Write its value
additively as \(z\), with
\((1+\rho+\rho^2)z=0\). Since three is invertible, set
\[
 b=\frac{(\rho+2\rho^2)z}{3}.
 \]
The identity
\((\rho-1)(\rho+2\rho^2)=3-(1+\rho+\rho^2)\) gives
\((\rho-1)b=z\). Thus the remaining cocycle is a coboundary.
For an actual pro-\(p\) quotient, apply this to all finite stable
quotients; the compact sets of conjugators have the
finite-intersection property, giving the same conclusion.

Every actual equivariant finite \(p\)-quotient of \(K'_3\)
therefore kills both classes. This excludes only that specified
transporter-only test. It neither equates the two actual compact
classes nor excludes relative-three or coupled pentagon targets.
\end{proof}

The complete predecessor proofs and original-source audit are retained
in \path{work_DE/original_r_frame_cohomology_gate33.md},
\path{work_DE/derived_kernel_cyclic_cohomology33.md}, and
\path{work_DE/original_CM_source_binding33.md}.
No finite census or abstract infinite-index completion is used.
